\begin{tikzpicture}[
  x=1mm,y=1mm,font=\figurefont\footnotesize,
  target/.style={circle,draw=FigureYellow,fill=FigureOutput!10,
    minimum size=8mm,inner sep=1pt},
  hidden/.style={rectangle,draw=FigureInk,fill=FigureModel!12,
    minimum size=8mm,inner sep=1pt},
  edge/.style={draw=FigureInk,line width=0.65pt},
  fill edge/.style={draw=FigureRed,dashed,line width=0.7pt},
  flow/.style={-{Stealth[length=2mm]},draw=FigureInk,line width=0.7pt}
]
  \node[text=FigureInk] at (20,23) {原始交互图};
  \node[text=FigureInk] at (82,23) {消去中心之后};
  \node[hidden] (c) at (20,0) {$X_0$};
  \node[target] (l1) at (5,13) {$X_1$};
  \node[target] (l2) at (35,13) {$X_2$};
  \node[target] (l3) at (5,-13) {$X_3$};
  \node[target] (l4) at (35,-13) {$X_4$};
  \foreach \j in {1,2,3,4}
    \draw[edge] (c)--(l\j);
  \node[target] (r1) at (67,13) {$X_1$};
  \node[target] (r2) at (97,13) {$X_2$};
  \node[target] (r3) at (67,-13) {$X_3$};
  \node[target] (r4) at (97,-13) {$X_4$};
  \draw[fill edge] (r1)--(r2);
  \draw[fill edge] (r1)--(r3);
  \draw[fill edge] (r1)--(r4);
  \draw[fill edge] (r2)--(r3);
  \draw[fill edge] (r2)--(r4);
  \draw[fill edge] (r3)--(r4);
  \draw[flow] (43,0)--node[above] {$\sum_{x_0}$}(59,0);
  \node[align=center,text=FigureMuted] at (20,-27)
    {方形：求和变量\\圆形：MAP目标变量};
  \node[align=center,text=FigureMuted] at (82,-27)
    {虚线：新增填充边\\中间因子含全部四个叶节点};
\end{tikzpicture}
